\IfFileExists{stacks-project.cls}{%
\documentclass{stacks-project}
}{%
\documentclass{amsart}
}

% For dealing with references we use the comment environment
\usepackage{verbatim}
\newenvironment{reference}{\comment}{\endcomment}
%\newenvironment{reference}{}{}
\newenvironment{slogan}{\comment}{\endcomment}
\newenvironment{history}{\comment}{\endcomment}

% For commutative diagrams we use Xy-pic
\usepackage[all]{xy}

% We use 2cell for 2-commutative diagrams.
\xyoption{2cell}
\UseAllTwocells

% We use multicol for the list of chapters between chapters
\usepackage{multicol}

% This is generally recommended for better output
\usepackage{lmodern}
\usepackage[T1]{fontenc}

% For cross-file-references
\usepackage{xr-hyper}

% Package for hypertext links:
\usepackage{hyperref}

% For any local file, say "hello.tex" you want to link to please
% use \externaldocument[hello-]{hello}
\externaldocument[introduction-]{introduction}
\externaldocument[conventions-]{conventions}
\externaldocument[sets-]{sets}
\externaldocument[categories-]{categories}
\externaldocument[topology-]{topology}
\externaldocument[sheaves-]{sheaves}
\externaldocument[sites-]{sites}
\externaldocument[stacks-]{stacks}
\externaldocument[fields-]{fields}
\externaldocument[algebra-]{algebra}
\externaldocument[brauer-]{brauer}
\externaldocument[homology-]{homology}
\externaldocument[derived-]{derived}
\externaldocument[simplicial-]{simplicial}
\externaldocument[more-algebra-]{more-algebra}
\externaldocument[smoothing-]{smoothing}
\externaldocument[modules-]{modules}
\externaldocument[sites-modules-]{sites-modules}
\externaldocument[injectives-]{injectives}
\externaldocument[cohomology-]{cohomology}
\externaldocument[sites-cohomology-]{sites-cohomology}
\externaldocument[dga-]{dga}
\externaldocument[dpa-]{dpa}
\externaldocument[sdga-]{sdga}
\externaldocument[hypercovering-]{hypercovering}
\externaldocument[schemes-]{schemes}
\externaldocument[constructions-]{constructions}
\externaldocument[properties-]{properties}
\externaldocument[morphisms-]{morphisms}
\externaldocument[coherent-]{coherent}
\externaldocument[divisors-]{divisors}
\externaldocument[limits-]{limits}
\externaldocument[varieties-]{varieties}
\externaldocument[topologies-]{topologies}
\externaldocument[descent-]{descent}
\externaldocument[perfect-]{perfect}
\externaldocument[more-morphisms-]{more-morphisms}
\externaldocument[flat-]{flat}
\externaldocument[groupoids-]{groupoids}
\externaldocument[more-groupoids-]{more-groupoids}
\externaldocument[etale-]{etale}
\externaldocument[chow-]{chow}
\externaldocument[intersection-]{intersection}
\externaldocument[pic-]{pic}
\externaldocument[weil-]{weil}
\externaldocument[adequate-]{adequate}
\externaldocument[dualizing-]{dualizing}
\externaldocument[duality-]{duality}
\externaldocument[discriminant-]{discriminant}
\externaldocument[derham-]{derham}
\externaldocument[local-cohomology-]{local-cohomology}
\externaldocument[algebraization-]{algebraization}
\externaldocument[curves-]{curves}
\externaldocument[resolve-]{resolve}
\externaldocument[models-]{models}
\externaldocument[functors-]{functors}
\externaldocument[equiv-]{equiv}
\externaldocument[pione-]{pione}
\externaldocument[etale-cohomology-]{etale-cohomology}
\externaldocument[proetale-]{proetale}
\externaldocument[relative-cycles-]{relative-cycles}
\externaldocument[more-etale-]{more-etale}
\externaldocument[trace-]{trace}
\externaldocument[crystalline-]{crystalline}
\externaldocument[spaces-]{spaces}
\externaldocument[spaces-properties-]{spaces-properties}
\externaldocument[spaces-morphisms-]{spaces-morphisms}
\externaldocument[decent-spaces-]{decent-spaces}
\externaldocument[spaces-cohomology-]{spaces-cohomology}
\externaldocument[spaces-limits-]{spaces-limits}
\externaldocument[spaces-divisors-]{spaces-divisors}
\externaldocument[spaces-over-fields-]{spaces-over-fields}
\externaldocument[spaces-topologies-]{spaces-topologies}
\externaldocument[spaces-descent-]{spaces-descent}
\externaldocument[spaces-perfect-]{spaces-perfect}
\externaldocument[spaces-more-morphisms-]{spaces-more-morphisms}
\externaldocument[spaces-flat-]{spaces-flat}
\externaldocument[spaces-groupoids-]{spaces-groupoids}
\externaldocument[spaces-more-groupoids-]{spaces-more-groupoids}
\externaldocument[bootstrap-]{bootstrap}
\externaldocument[spaces-pushouts-]{spaces-pushouts}
\externaldocument[spaces-chow-]{spaces-chow}
\externaldocument[groupoids-quotients-]{groupoids-quotients}
\externaldocument[spaces-more-cohomology-]{spaces-more-cohomology}
\externaldocument[spaces-simplicial-]{spaces-simplicial}
\externaldocument[spaces-duality-]{spaces-duality}
\externaldocument[formal-spaces-]{formal-spaces}
\externaldocument[restricted-]{restricted}
\externaldocument[spaces-resolve-]{spaces-resolve}
\externaldocument[formal-defos-]{formal-defos}
\externaldocument[defos-]{defos}
\externaldocument[cotangent-]{cotangent}
\externaldocument[examples-defos-]{examples-defos}
\externaldocument[algebraic-]{algebraic}
\externaldocument[examples-stacks-]{examples-stacks}
\externaldocument[stacks-sheaves-]{stacks-sheaves}
\externaldocument[criteria-]{criteria}
\externaldocument[artin-]{artin}
\externaldocument[quot-]{quot}
\externaldocument[stacks-properties-]{stacks-properties}
\externaldocument[stacks-morphisms-]{stacks-morphisms}
\externaldocument[stacks-limits-]{stacks-limits}
\externaldocument[stacks-cohomology-]{stacks-cohomology}
\externaldocument[stacks-perfect-]{stacks-perfect}
\externaldocument[stacks-introduction-]{stacks-introduction}
\externaldocument[stacks-more-morphisms-]{stacks-more-morphisms}
\externaldocument[stacks-geometry-]{stacks-geometry}
\externaldocument[moduli-]{moduli}
\externaldocument[moduli-curves-]{moduli-curves}
\externaldocument[examples-]{examples}
\externaldocument[exercises-]{exercises}
\externaldocument[guide-]{guide}
\externaldocument[desirables-]{desirables}
\externaldocument[coding-]{coding}
\externaldocument[obsolete-]{obsolete}
\externaldocument[fdl-]{fdl}
\externaldocument[index-]{index}

% Theorem environments.
%
\theoremstyle{plain}
\newtheorem{theorem}[subsection]{Theorem}
\newtheorem{proposition}[subsection]{Proposition}
\newtheorem{lemma}[subsection]{Lemma}

\theoremstyle{definition}
\newtheorem{definition}[subsection]{Definition}
\newtheorem{example}[subsection]{Example}
\newtheorem{exercise}[subsection]{Exercise}
\newtheorem{situation}[subsection]{Situation}

\theoremstyle{remark}
\newtheorem{remark}[subsection]{Remark}
\newtheorem{remarks}[subsection]{Remarks}

\numberwithin{equation}{subsection}

% Macros
%
\def\lim{\mathop{\mathrm{lim}}\nolimits}
\def\colim{\mathop{\mathrm{colim}}\nolimits}
\def\Spec{\mathop{\mathrm{Spec}}}
\def\Hom{\mathop{\mathrm{Hom}}\nolimits}
\def\Ext{\mathop{\mathrm{Ext}}\nolimits}
\def\SheafHom{\mathop{\mathcal{H}\!\mathit{om}}\nolimits}
\def\SheafExt{\mathop{\mathcal{E}\!\mathit{xt}}\nolimits}
\def\Sch{\mathit{Sch}}
\def\Mor{\mathop{\mathrm{Mor}}\nolimits}
\def\Ob{\mathop{\mathrm{Ob}}\nolimits}
\def\Sh{\mathop{\mathit{Sh}}\nolimits}
\def\NL{\mathop{N\!L}\nolimits}
\def\CH{\mathop{\mathrm{CH}}\nolimits}
\def\proetale{{pro\text{-}\acute{e}tale}}
\def\etale{{\acute{e}tale}}
\def\QCoh{\mathit{QCoh}}
\def\Ker{\mathop{\mathrm{Ker}}}
\def\Im{\mathop{\mathrm{Im}}}
\def\Coker{\mathop{\mathrm{Coker}}}
\def\Coim{\mathop{\mathrm{Coim}}}

% Boxtimes
%
\DeclareMathSymbol{\boxtimes}{\mathbin}{AMSa}{"02}

%
% Macros for moduli stacks/spaces
%
\def\QCohstack{\mathcal{QC}\!\mathit{oh}}
\def\Cohstack{\mathcal{C}\!\mathit{oh}}
\def\Spacesstack{\mathcal{S}\!\mathit{paces}}
\def\Quotfunctor{\mathrm{Quot}}
\def\Hilbfunctor{\mathrm{Hilb}}
\def\Curvesstack{\mathcal{C}\!\mathit{urves}}
\def\Polarizedstack{\mathcal{P}\!\mathit{olarized}}
\def\Complexesstack{\mathcal{C}\!\mathit{omplexes}}
% \Pic is the operator that assigns to X its picard group, usage \Pic(X)
% \Picardstack_{X/B} denotes the Picard stack of X over B
% \Picardfunctor_{X/B} denotes the Picard functor of X over B
\def\Pic{\mathop{\mathrm{Pic}}\nolimits}
\def\Picardstack{\mathcal{P}\!\mathit{ic}}
\def\Picardfunctor{\mathrm{Pic}}
\def\Deformationcategory{\mathcal{D}\!\mathit{ef}}
\usepackage{cleveref}
\crefname{originallemma}{original lemma}{original lemmas}
\Crefname{originallemma}{Original lemma}{Original lemmas}
\makeatletter
\newlabel{lemma-get-cycles}{{0H59}{source}{}{sourcepointer-1}{}}
\newlabel{lemma-get-cycles@cref}{{[originallemma][0][]0H59}{[1][0][]source}{}{}{}}
\makeatother
\crefname{originalequation}{original equation}{original equations}
\Crefname{originalequation}{Original equation}{Original equations}
\makeatletter
\newlabel{equation-cycle-classes-general}{{0H6V}{source}{}{sourcepointer-2}{}}
\newlabel{equation-cycle-classes-general@cref}{{[originalequation][0][]0H6V}{[1][0][]source}{}{}{}}
\makeatother
\crefname{originalproposition}{original proposition}{original propositions}
\Crefname{originalproposition}{Original proposition}{Original propositions}
\makeatletter
\newlabel{limits-proposition-approximate}{{01ZA}{source}{}{sourcepointer-3}{}}
\newlabel{limits-proposition-approximate@cref}{{[originalproposition][0][]01ZA}{[1][0][]source}{}{}{}}
\makeatother
\crefname{originallemma}{original lemma}{original lemmas}
\Crefname{originallemma}{Original lemma}{Original lemmas}
\makeatletter
\newlabel{limits-lemma-descend-finite-presentation}{{01ZM}{source}{}{sourcepointer-4}{}}
\newlabel{limits-lemma-descend-finite-presentation@cref}{{[originallemma][0][]01ZM}{[1][0][]source}{}{}{}}
\makeatother
\crefname{originallemma}{original lemma}{original lemmas}
\Crefname{originallemma}{Original lemma}{Original lemmas}
\makeatletter
\newlabel{limits-lemma-descend-closed-immersion-finite-presentation}{{01ZP}{source}{}{sourcepointer-5}{}}
\newlabel{limits-lemma-descend-closed-immersion-finite-presentation@cref}{{[originallemma][0][]01ZP}{[1][0][]source}{}{}{}}
\makeatother
\crefname{originallemma}{original lemma}{original lemmas}
\Crefname{originallemma}{Original lemma}{Original lemmas}
\makeatletter
\newlabel{limits-lemma-descend-flat-finite-presentation}{{04AI}{source}{}{sourcepointer-6}{}}
\newlabel{limits-lemma-descend-flat-finite-presentation@cref}{{[originallemma][0][]04AI}{[1][0][]source}{}{}{}}
\makeatother
\crefname{originallemma}{original lemma}{original lemmas}
\Crefname{originallemma}{Original lemma}{Original lemmas}
\makeatletter
\newlabel{limits-lemma-descend-surjective}{{07RR}{source}{}{sourcepointer-7}{}}
\newlabel{limits-lemma-descend-surjective@cref}{{[originallemma][0][]07RR}{[1][0][]source}{}{}{}}
\makeatother
\crefname{originallemma}{original lemma}{original lemmas}
\Crefname{originallemma}{Original lemma}{Original lemmas}
\makeatletter
\newlabel{limits-lemma-eventually-proper}{{081F}{source}{}{sourcepointer-8}{}}
\newlabel{limits-lemma-eventually-proper@cref}{{[originallemma][0][]081F}{[1][0][]source}{}{}{}}
\makeatother
\crefname{originallemma}{original lemma}{original lemmas}
\Crefname{originallemma}{Original lemma}{Original lemmas}
\makeatletter
\newlabel{limits-lemma-limit-dimension}{{05M5}{source}{}{sourcepointer-9}{}}
\newlabel{limits-lemma-limit-dimension@cref}{{[originallemma][0][]05M5}{[1][0][]source}{}{}{}}
\makeatother
\crefname{originallemma}{original lemma}{original lemmas}
\Crefname{originallemma}{Original lemma}{Original lemmas}
\makeatletter
\newlabel{more-morphisms-lemma-descend-cd}{{0GTM}{source}{}{sourcepointer-10}{}}
\newlabel{more-morphisms-lemma-descend-cd@cref}{{[originallemma][0][]0GTM}{[1][0][]source}{}{}{}}
\makeatother
\crefname{originallemma}{original lemma}{original lemmas}
\Crefname{originallemma}{Original lemma}{Original lemmas}
\makeatletter
\newlabel{lemma-relative-cycles-equal}{{0H5C}{source}{}{sourcepointer-11}{}}
\newlabel{lemma-relative-cycles-equal@cref}{{[originallemma][0][]0H5C}{[1][0][]source}{}{}{}}
\makeatother
\crefname{originallemma}{original lemma}{original lemmas}
\Crefname{originallemma}{Original lemma}{Original lemmas}
\makeatletter
\newlabel{limits-lemma-limit-contained-in-constructible}{{05F4}{source}{}{sourcepointer-12}{}}
\newlabel{limits-lemma-limit-contained-in-constructible@cref}{{[originallemma][0][]05F4}{[1][0][]source}{}{}{}}
\makeatother
\crefname{originallemma}{original lemma}{original lemmas}
\Crefname{originallemma}{Original lemma}{Original lemmas}
\makeatletter
\newlabel{lemma-descend-family}{{0H4X}{source}{}{sourcepointer-13}{}}
\newlabel{lemma-descend-family@cref}{{[originallemma][0][]0H4X}{[1][0][]source}{}{}{}}
\makeatother
\crefname{originallemma}{original lemma}{original lemmas}
\Crefname{originallemma}{Original lemma}{Original lemmas}
\makeatletter
\newlabel{lemma-relative-cycles-h-descent}{{0H52}{source}{}{sourcepointer-14}{}}
\newlabel{lemma-relative-cycles-h-descent@cref}{{[originallemma][0][]0H52}{[1][0][]source}{}{}{}}
\makeatother
\title{Descent of a family of relative cycles to a Noetherian model, criterion via completely decomposed proper morphisms}
\author{The Stacks Project Authors}
\date{stacks-project:relative-cycles.tex:local-snapshot-20261012; GNU FDL 1.2 or later}
\begin{document}
\maketitle
\setcounter{subsection}{1}
\noindent\textit{Editorial note.} All delivered mathematical text is the author's own, copied from the pinned Stacks Project chapter file relative-cycles.tex; no statement or proof step was rewritten, reordered or removed. The counted span is the author environment and its own proof environment. References inside the span are handed to the assembler as navigation stubs that carry the canonical tag of the cited tag only; whether each of them is load-bearing is recorded on the stub itself.\par\medskip
\noindent\textit{Allowed standard prerequisites (editorial).}\begin{itemize}
\item Existence of a completely decomposed proper morphism g\_0 : S'\_0 to S\_0 putting a given family on a Noetherian base into the image of the cycle-classes construction (relative-cycles chapter, Lemma lemma-get-cycles, canonical tag 0H59). cited at proof lines 3179-3181. Direction (1) to (2): produces the witness for the Noetherian model, which base changes to S.
\item The cycle-classes construction whose image is being characterized (relative-cycles chapter, Equation equation-cycle-classes-general, canonical tag 0H6V). cited at proof lines 3173, 3181 and 3188. The map z(X/S, r) is characterized against; being in its image is condition (2).
\item Approximation: a quasi-compact quasi-separated scheme is a directed limit of finite type schemes over Z with affine transitions (limits chapter, Proposition limits-proposition-approximate, canonical tag 01ZA). cited at proof line 3197. Sets up writing S = lim S\_i for the descent of the data in direction (2) to (1).
\item Descent of finite presentation along limits (limits chapter, Lemma limits-lemma-descend-finite-presentation, canonical tag 01ZM). cited at proof line 3208. One of the six Limits inputs descending the morphism and subscheme data to some finite stage.
\item Descent of closed immersions of finite presentation along limits (limits chapter, Lemma limits-lemma-descend-closed-immersion-finite-presentation, canonical tag 01ZP). cited at proof line 3209. Descends the closed subschemes Z\_a to Z\_\{ai\}.
\item Descent of flatness of finite presentation along limits (limits chapter, Lemma limits-lemma-descend-flat-finite-presentation, canonical tag 04AI). cited at proof line 3210. Preserves the flatness of Z\_\{ai\} over S'\_i.
\item Descent of surjectivity along limits (limits chapter, Lemma limits-lemma-descend-surjective, canonical tag 07RR). cited at proof line 3211. Preserves complete decomposition of the descended g\_i.
\item Eventual properness along limits (limits chapter, Lemma limits-lemma-eventually-proper, canonical tag 081F). cited at proof line 3212. Makes the descended g\_i proper at some finite stage.
\item Dimension bounds along limits (limits chapter, Lemma limits-lemma-limit-dimension, canonical tag 05M5). cited at proof line 3213. Preserves the relative dimension bound of the descended subschemes.
\item Descent of completely decomposed morphisms (more-morphisms chapter, Lemma more-morphisms-lemma-descend-cd, canonical tag 0GTM). cited at proof line 3216. Makes the descended g\_i completely decomposed.
\item Closedness of the locus where two relative cycles on the fibres agree (relative-cycles chapter, Lemma lemma-relative-cycles-equal, canonical tag 0H5C). cited at proof line 3232. The equal-locus E is closed, enabling the constructible-limit argument.
\item A morphism from a limit scheme landing in a constructible subset factors through some finite stage (limits chapter, Lemma limits-lemma-limit-contained-in-constructible, canonical tag 05F4). cited at proof lines 3234-3236. Promotes the inclusion of S'' into E to a finite stage i', making the two pullbacks literally equal.
\item Unique descent of a family of cycles on fibres along a completely decomposed proper morphism (relative-cycles chapter, Lemma lemma-descend-family, canonical tag 0H4X). cited at proof lines 3239-3242. Produces the unique family alpha\_i on X\_i/S\_i with g\_i\textasciicircum{}* alpha\_i = alpha'\_i.
\item The h-descent property placing descended families into the relative cycles group z(X\_i/S\_i, r) (relative-cycles chapter, Lemma lemma-relative-cycles-h-descent, canonical tag 0H52). cited at proof lines 3244-3245. Verifies alpha\_i is an honest relative r-cycle, completing condition (1).
\end{itemize}\medskip
\ifdefined\hypertarget\hypertarget{sourcepointer-1}{}\fi
\noindent\textit{Source reference.} Original lemma 0H59: sources/stacks-chow/tags.txt, source lines 21379--21379. Cited at proof lines 3179-3181: direction (1) to (2) rests on it.\par\medskip
\ifdefined\hypertarget\hypertarget{sourcepointer-2}{}\fi
\noindent\textit{Source reference.} Original equation 0H6V: sources/stacks-chow/tags.txt, source lines 21435--21435. Cited at proof lines 3173, 3181, 3188: the image condition (2) is formulated against it.\par\medskip
\ifdefined\hypertarget\hypertarget{sourcepointer-3}{}\fi
\noindent\textit{Source reference.} Original proposition 01ZA: sources/stacks-chow/tags.txt, source lines 2585--2585. Cited at proof line 3197: writes S as a limit of Noetherian finite type stages.\par\medskip
\ifdefined\hypertarget\hypertarget{sourcepointer-4}{}\fi
\noindent\textit{Source reference.} Original lemma 01ZM: sources/stacks-chow/tags.txt, source lines 2599--2599. Cited at proof line 3208: descends the morphism data to a finite stage.\par\medskip
\ifdefined\hypertarget\hypertarget{sourcepointer-5}{}\fi
\noindent\textit{Source reference.} Original lemma 01ZP: sources/stacks-chow/tags.txt, source lines 2602--2602. Cited at proof line 3209: descends the closed subschemes Z\_a.\par\medskip
\ifdefined\hypertarget\hypertarget{sourcepointer-6}{}\fi
\noindent\textit{Source reference.} Original lemma 04AI: sources/stacks-chow/tags.txt, source lines 5617--5617. Cited at proof line 3210: keeps the descended subschemes flat.\par\medskip
\ifdefined\hypertarget\hypertarget{sourcepointer-7}{}\fi
\noindent\textit{Source reference.} Original lemma 07RR: sources/stacks-chow/tags.txt, source lines 9876--9876. Cited at proof line 3211: keeps the descended cover surjective.\par\medskip
\ifdefined\hypertarget\hypertarget{sourcepointer-8}{}\fi
\noindent\textit{Source reference.} Original lemma 081F: sources/stacks-chow/tags.txt, source lines 10215--10215. Cited at proof line 3212: makes the descended g\_i proper.\par\medskip
\ifdefined\hypertarget\hypertarget{sourcepointer-9}{}\fi
\noindent\textit{Source reference.} Original lemma 05M5: sources/stacks-chow/tags.txt, source lines 7259--7259. Cited at proof line 3213: keeps the relative dimension bound of the descended subschemes.\par\medskip
\ifdefined\hypertarget\hypertarget{sourcepointer-10}{}\fi
\noindent\textit{Source reference.} Original lemma 0GTM: sources/stacks-chow/tags.txt, source lines 20972--20972. Cited at proof line 3216: makes the descended g\_i completely decomposed.\par\medskip
\ifdefined\hypertarget\hypertarget{sourcepointer-11}{}\fi
\noindent\textit{Source reference.} Original lemma 0H5C: sources/stacks-chow/tags.txt, source lines 21382--21382. Cited at proof line 3232: makes the equal locus E closed.\par\medskip
\ifdefined\hypertarget\hypertarget{sourcepointer-12}{}\fi
\noindent\textit{Source reference.} Original lemma 05F4: sources/stacks-chow/tags.txt, source lines 7013--7013. Cited at proof lines 3234-3236: promotes the equal-locus inclusion to a finite stage.\par\medskip
\ifdefined\hypertarget\hypertarget{sourcepointer-13}{}\fi
\noindent\textit{Source reference.} Original lemma 0H4X: sources/stacks-chow/tags.txt, source lines 21367--21367. Cited at proof lines 3239-3242: descends alpha'\_i uniquely to alpha\_i.\par\medskip
\ifdefined\hypertarget\hypertarget{sourcepointer-14}{}\fi
\noindent\textit{Source reference.} Original lemma 0H52: sources/stacks-chow/tags.txt, source lines 21372--21372. Cited at proof lines 3244-3245: verifies alpha\_i lies in z(X\_i/S\_i, r).\par\medskip
\begin{lemma}
\label{lemma-relative-r-cycle-general}
Let $S$ be a quasi-compact and quasi-separated scheme.
Let $f : X \to S$ be a morphism of finite presentation.
Let $r \geq 0$ and let $\alpha$ be a family of $r$-cycles on fibres of $X/S$.
The following are equivalent
\begin{enumerate}
\item there exists a cartesian diagram
$$
\xymatrix{
X \ar[r] \ar[d] & X_0 \ar[d] \\
S \ar[r] & S_0
}
$$
where $X_0 \to S_0$ is a finite type morphism of Noetherian schemes
and $\alpha_0 \in z(X_0/S_0, r)$ such that $\alpha$ is the base change
of $\alpha_0$ by $S \to S_0$
\item there exists a completely decomposed proper morphism $g : S' \to S$
of finite presentation such that $g^*\alpha$ is in the image of
(\ref{equation-cycle-classes-general}).
\end{enumerate}
\end{lemma}

\begin{proof}
Let a diagram and $\alpha_0 \in z(X_0/S_0, r)$ as in (1) be given.
By Lemma \ref{lemma-get-cycles} there exists a completely decomposed proper
morphism $g_0 : S'_0 \to S_0$ such that $g_0^*\alpha_0$ is in
the image of (\ref{equation-cycle-classes-general}). Namely, since
$S'_0$ is Noetherian, every closed subscheme of $S'_0 \times_{S_0} X_0$
is of finite presentation over $S'_0$. Setting $S' = S \times_{S_0} S'_0$
and using base change by $S' \to S'_0$ we see that (2) holds.

\medskip\noindent
Conversely, assume that (2) holds. Choose a completely decomposed proper
morphism $g : S' \to S$ of finite presentation such that $g^*\alpha$ is in the
image of (\ref{equation-cycle-classes-general}). Set $X' = S' \times_S X$.
Write $g^*\alpha = \sum n_a [Z_a/X'/S']_r$ for some $Z_a \subset X'$
closed subscheme flat, of finite presentation, and of relative dimension
$\leq r$ over $S'$.

\medskip\noindent
Write $S = \lim S_i$ as a directed limit with affine transition
morphisms with $S_i$ of finite type over $\mathbf{Z}$,
see Limits, Proposition \ref{limits-proposition-approximate}.
We can find an $i$ large enough such that there exist
\begin{enumerate}
\item a completely decomposed proper morphism $g_i : S'_i \to S_i$
whose base change to $S$ is $g : S' \to S$,
\item setting $X'_i = S'_i \times_{S_i} X_i$
closed subschemes $Z_{ai} \subset X'_i$ flat and of
relative dimension $\leq r$ over $S'_i$ whose base change to $S'$
is $Z_a$.
\end{enumerate}
To do this one uses Limits, Lemmas
\ref{limits-lemma-descend-finite-presentation},
\ref{limits-lemma-descend-closed-immersion-finite-presentation},
\ref{limits-lemma-descend-flat-finite-presentation},
\ref{limits-lemma-descend-surjective},
\ref{limits-lemma-eventually-proper}, and
\ref{limits-lemma-limit-dimension}
and
and More on Morphisms, Lemma
\ref{more-morphisms-lemma-descend-cd}.
Consider $\alpha'_i = \sum n_a [Z_{ai}/X'_i/S_i]_r \in z(X'_i/S'_i, r)$.
The base change of $\alpha'_i$ to a family of $r$-cycles on fibres
of $X'/S'$ agrees with the base change $g^*\alpha$ by construction.

\medskip\noindent
Set $S''_i = S'_i \times_{S_i} S'_i$ and $X''_i = S''_i \times_{S_i} X_i$
and set $S'' = S' \times_S S'$ and $X'' = S'' \times_S X$.
We denote $\text{pr}_1, \text{pr}_2 : S'' \to S'$ and
$\text{pr}_1, \text{pr}_2 : S''_i \to S'_i$ the projections.
The relative $r$-cycles $\text{pr}_1^*\alpha'_i$ and $\text{pr}_1^*\alpha'_i$
on $X''_i/S''_i$ base change to the same family of $r$-cycles
on fibres of $X''/S''$ because
$\text{pr}_1^*g^*\alpha = \text{pr}_1^*g^*\alpha$.
Hence the morphism $S'' \to S''_i$ maps into $E =
\{s \in S''_i : (\text{pr}_1^*\alpha'_i)_s = (\text{pr}_1^*\alpha'_i)_s\}$.
By Lemma \ref{lemma-relative-cycles-equal} this is a closed subset.
Since $S'' = \lim_{i' \geq i} S''_{i'}$ we see from
Limits, Lemma
\ref{limits-lemma-limit-contained-in-constructible}
that for some $i' \geq i$ the morphism $S''_{i'} \to S''_i$
maps into $E$. Therefore, after replacing $i$ by $i'$, we may assume that
$\text{pr}_1^*\alpha'_i = \text{pr}_1^*\alpha'_i$.
By Lemma \ref{lemma-descend-family}
we obtain a unique family $\alpha_i$
of $r$-cycles on fibres of $X_i/S_i$
with $g_i^*\alpha_i = \alpha'_i$ (this uses that $S'_i \to S_i$
is completely decomposed).
By Lemma \ref{lemma-relative-cycles-h-descent}
we see that $\alpha_i \in z(X_i/S_i, r)$.
The uniqueness in Lemma \ref{lemma-descend-family} implies that the
base change of $\alpha_i$ is $\alpha$ and we see (1) holds.
\end{proof}

\end{document}
